\documentclass[10pt,a4paper]{amsart}
\usepackage[margin=19mm]{geometry}
\usepackage{amsmath,amssymb,mathtools}
\usepackage[scr=rsfs,cal=euler]{mathalfa}
\usepackage{xfrac}
\usepackage[unicode,hidelinks,bookmarks=false]{hyperref}
\newcommand{\ie}{i.e.\ }
\newcommand{\cf}{cf.\ }
\newcommand{\resp}{respectively\ }
\newcommand{\Subsection}{\subsection}
\providecommand{\nobreakdash}{\nobreak\hbox{-}}
\setlength{\emergencystretch}{2em}
\multlinegap0pt
\usepackage{amssymb}
\usepackage{mathtools}
\usepackage[shortlabels]{enumitem}
\newtheorem{lemma}{Lemma}
\newtheorem{proposition}{Proposition}
\newcommand{\R}{\mathbb{R}}
\newcommand{\cT}{\mathcal{R}}
\newcommand{\dd}{\,\mathrm{d}}
\newcommand{\eps}{\varepsilon}
\newcommand{\norm}[1]{\left\lVert #1\right\rVert}
\newcommand{\nupar}{\nu}
\title{Regularized Subjective-Surface Flow with Monotone Reaction: Global Classical Well-Posedness and Stability}
\author{Markjoe O. Uba}
\date{Source: arXiv:2608.22015v1 (CC BY 4.0)}
\begin{document}
\maketitle

\noindent\emph{Source and licence.} Regularized Subjective-Surface Flow with Monotone Reaction: Global Classical Well-Posedness and Stability. Version arXiv:2608.22015v1 (CC BY 4.0), distributed by arXiv under the Creative Commons Attribution 4.0 International licence, http://creativecommons.org/licenses/by/4.0/, as recorded in the arXiv licence metadata for this exact version and shown on the abstract page https://arxiv.org/abs/2608.22015v1. Full licence notice: \texttt{licenses/arxiv-2608.22015-CC-BY-4.0.txt}.\par\smallskip

\noindent\textit{Editorial scope.} Counted unit: the article's proposition prop:finite-gradient (source0.tex lines 842--853). The article's complete original proof of it is delivered with it (source0.tex lines 855--1074, one \texttt{proof} environment of 220 lines). No mathematics is written, repaired or re-derived: the statement and the proof are the article's own lines, in the article's own notation, and no retained line is substituted or re-flowed.  Retained uncounted as the source-local context the proof needs: lines 208--219; lines 393--399; lines 450--460; lines 713--731.  Nothing was omitted from any retained span, so the package carries no omission record.  No retained line contains a citation, so the wrapper carries no bibliography and no bibliography entry.  No retained line contains a figure environment, an image-inclusion command or drawing code, so no image is delivered and the package declares no asset dependency.  Editorial formatting only, in addition: an AMS \texttt{amsart} preamble that restates the article's own declarations and macros used by the retained lines, namely the environments \texttt{lemma}, \texttt{proposition} on separate counters, together with the standard packages those lines need.  The retained environments print in this document as Lemma 1, Proposition 1 in order of appearance, whereas the article prints the same items with the numbering of its own full document.  The complete native source of the article as distributed in the arXiv e-print for this exact version is bundled unchanged as \texttt{sources/208294/source0.tex}.
\begin{equation}\label{eq:main-pde}
\boxed{
\begin{cases}
\partial_t u
=\nupar\Delta u+
A_\eps(\nabla u)\operatorname{div}\!\left(G(x)\dfrac{\nabla u}{A_\eps(\nabla u)}\right)
-\mu\cT_\eta(x,u),
& (t,x)\in(0,\infty)\times\Omega,\\[2mm]
u=0,& (t,x)\in(0,\infty)\times\partial\Omega,\\[1mm]
u(0,x)=u^0(x),& x\in\Omega.
\end{cases}}
\end{equation}

\begin{lemma}[Uniform parabolicity and global coefficient bounds]\label{lem:ellipticity}
For every $x\in\overline\Omega$ and $p,\xi\in\R^d$,
\begin{equation}\label{eq:ellipticity}
        \nupar |\xi|^2\le a^{ij}(x,p)\xi_i\xi_j\le (\nupar+g^*)|\xi|^2.
\end{equation}
Moreover, for $m=1,2$, the derivatives $D_p^m a^{ij}(x,p)$ are bounded on $\overline\Omega\times\R^d$ by constants depending only on $d$, $m$, $\eps$, and $g^*$.  The $x$-derivatives up to order two are bounded by $\norm{G}_{C^{2}(\overline\Omega)}$, and the H\"older moduli used in the positive-time estimates are controlled by $\norm{G}_{C^{3+\alpha}(\overline\Omega)}$.  In particular, these coefficient bounds do not require an a priori bound on $|p|$.
\end{lemma}

\begin{lemma}[Comparison principle on a positive-time interval]\label{lem:comparison}
Let $0\le t_0<T$ and
\[
 w\in C([t_0,T];C^1(\overline\Omega))\cap C^{1,2}((t_0,T]\times\overline\Omega)
\]
satisfy
\[
 w_t-a^{ij}(t,x)w_{ij}-c^i(t,x)w_i+d(t,x)w\le0
\]
in $(t_0,T]\times\Omega$.  Assume that $a^{ij},c^i,d$ are continuous and bounded on $[t_0,T]\times\overline\Omega$, that $a^{ij}$ is uniformly parabolic, and that $d\ge0$.  If $w\le0$ at $t=t_0$ and on the spatial boundary, then $w\le0$ in $[t_0,T]\times\overline\Omega$.
\end{lemma}

\begin{lemma}[Interior regularity improvement for positive times]
\label{lem:interior-bootstrap}
Let $u$ be a positive-time classical solution of
\eqref{eq:main-pde}. For every pair of parabolic cylinders
\[
Q'\Subset Q\Subset (0,T]\times\Omega,
\]
there exists $\sigma\in(0,\alpha]$ such that
\[
D_xu\in C^{1+\sigma/2,\,2+\sigma}(Q').
\]
Equivalently,
\[
u\in C^{1+\sigma/2,\,3+\sigma}(Q').
\]
Consequently, all derivatives required to differentiate the equation once
with respect to the spatial variables exist and are H\"older continuous at
every positive-time interior point.
\end{lemma}

\begin{proposition}[Finite-time gradient estimate]\label{prop:finite-gradient}
Let $U$ be a positive-time classical solution of \eqref{eq:main-pde} on $[0,S]$ with $U(t,x)=0$ for $t>0$ and $x\in\partial\Omega$, initial datum $U^0\in C^{2+\alpha}(\overline\Omega)$ satisfying $U^0=0$ on $\partial\Omega$, $U(t)\to U^0$ in $C^1(\overline\Omega)$ as $t\downarrow0$, and $0\le U\le1$.  Then there exists a constant
\[
        C_{\nabla}=C_{\nabla}\bigl(S,\Omega,d,\eps,\nupar,g_*,g^*,
        \norm{G}_{C^{3+\alpha}},\norm{\Lambda}_{C^{1+\alpha}},
        \mu,q,\eta,H,\norm{U^0}_{C^{2+\alpha}}\bigr)
\]
such that
\[
        \norm{\nabla U}_{L^\infty((0,S)\times\Omega)}\le C_{\nabla}.
\]
\end{proposition}

\begin{proof}
Write the equation in nondivergence form as
\begin{equation}\label{eq:tailored-gradient-pde}
        U_t=a^{ij}(x,\nabla U)U_{ij}+b^i(x)U_i+F(x,U),
\end{equation}
where
\[
        b=\nabla G,
        \qquad
        F(x,r)=-\mu\Lambda(x)H_\eta(r-q),
\]
and
\[
        a^{ij}(x,p)=\nupar\delta_{ij}
        +G(x)\left(\delta_{ij}-\frac{p_i p_j}{\eps^2+|p|^2}\right).
\]
By Lemma~\ref{lem:ellipticity},
\[
        \nupar|\xi|^2\le a^{ij}(x,p)\xi_i\xi_j
        \le (\nupar+g^*)|\xi|^2.
\]
Moreover $D_xa^{ij}$ is bounded independently of $p$, $D_p a^{ij}$ is bounded because $\eps>0$, $b\in C^{2+\alpha}(\overline\Omega)$, and $F_x,F_r$ are bounded on $\overline\Omega\times[0,1]$, with
\[
        F_r(x,r)=-\mu\Lambda(x)H_\eta'(r-q)\le0.
\]

We first control the boundary gradient.  Since $\partial\Omega$ is $C^{4+\alpha}$, there are $\delta_0>0$ and $M_\rho>0$ such that the interior distance function
\[
        \rho(x)=\operatorname{dist}(x,\partial\Omega)
\]
is $C^2$ in the neighborhood of the boundary $\Omega_{\delta_0}=\{x\in\Omega:0<\rho(x)<\delta_0\}$ and
\[
        |D^2\rho|\le M_\rho \quad\text{there}.
\]
Let $n=\nabla\rho$ be the inward unit normal.  Choose $\kappa>0$ so large that
\[
        \kappa\nupar
        \ge
        1+d(\nupar+g^*)M_\rho+\norm{\nabla G}_{L^\infty(\Omega)}.
\]
For constants $M>0$ to be fixed, set
\[
        W(x)=M\bigl(1-e^{-\kappa\rho(x)}\bigr).
\]
Increasing $M$ if necessary, we may arrange that
\[
        W\ge1 \quad\text{on } \{\rho=\delta_0\},
        \qquad
        W\ge U^0 \quad\text{in } \Omega_{\delta_0}.
\]
The second inequality is possible because $U^0=0$ on $\partial\Omega$ and $U^0\in C^1(\overline\Omega)$; equivalently, the quotient
\[
        \frac{U^0(x)}{1-e^{-\kappa\rho(x)}}
\]
extends boundedly to the boundary.

In the neighborhood of the boundary,
\[
        \nabla W=M\kappa e^{-\kappa\rho}n,
        \qquad
        D^2W=-M\kappa^2 e^{-\kappa\rho} n\otimes n
        +M\kappa e^{-\kappa\rho}D^2\rho.
\]
Equivalently, if $\phi(s)=M(1-e^{-\kappa s})$, then $\phi'>0$ and $\phi''=-\kappa\phi'$, so
\[
\begin{aligned}
        &W_t-a^{ij}(x,\nabla W)W_{ij}-b(x)\cdot\nabla W
        +\mu\Lambda(x)H_\eta(W-q) \\
        &\qquad
        =-a^{ij}(x,\phi'n)\bigl(\phi'' n_i n_j+\phi'\rho_{ij}\bigr)
        -\phi'\nabla G\cdot n
        +\mu\Lambda(x)H_\eta(W-q) \\
        &\qquad
        \ge
        \phi'\bigl(\kappa\nupar-d(\nupar+g^*)M_\rho-\norm{\nabla G}_{L^\infty}\bigr)
        \ge0.
\end{aligned}
\]
Thus $W$ is a supersolution in the neighborhood of the boundary.  We compare it with $U$ by applying Lemma~\ref{lem:comparison} to the linearized difference $z=U-W$.  Indeed,
\[
\begin{aligned}
        &a^{ij}(x,\nabla U)U_{ij}-a^{ij}(x,\nabla W)W_{ij} \\
        &\qquad =a^{ij}(x,\nabla U)z_{ij}
        +\left(\int_0^1 \partial_{p_k}a^{ij}\bigl(x,\nabla W+s\nabla z\bigr)\,ds\right)W_{ij}z_k,
\end{aligned}
\]
and
\[
        \cT_\eta(x,U)-\cT_\eta(x,W)
        =\left(\int_0^1 \partial_r\cT_\eta(x,W+s z)\,ds\right)z .
\]
Equivalently, $z$ satisfies an inequality of the form
\[
        z_t-\tilde a^{ij}z_{ij}-\tilde c^k z_k+\tilde d z\le0,
\]
where $\tilde a^{ij}=a^{ij}(x,\nabla U)$ is uniformly elliptic and $\tilde c^k$ is bounded in the boundary neighborhood.  Indeed, this bound uses only the global bound for $D_pa$ and the fixed $C^2$ norm of $W$, and hence does not depend on an a priori bound for $\nabla U$.  Moreover,
\[
        \tilde d
        =\mu\int_0^1\Lambda(x)H_\eta'(W+s z-q)\,\dd s\ge0.
\]
This explicit sign follows exactly from the monotonicity of $r\mapsto H_\eta(r-q)$.  To avoid using the equation at the initial corner $t=0$, fix $\tau>0$ and set
\[
 \varepsilon_\tau=\bigl\|(U(\tau)-W)_+\bigr\|_{L^\infty(\Omega_{\delta_0})}.
\]
Because $U(t)\to U^0$ in $C^1(\overline\Omega)$ and $U^0\le W$, one has $\varepsilon_\tau\to0$.  Since adding a nonnegative constant leaves the principal part unchanged and increases the nondecreasing reaction, $W+\varepsilon_\tau$ is again a supersolution.  On the parabolic boundary of $[\tau,S]\times\Omega_{\delta_0}$ we have
$U(\tau)\le W+\varepsilon_\tau$, $U=0\le W+\varepsilon_\tau$ on $\{\rho=0\}$, and $U\le1\le W+\varepsilon_\tau$ on $\{\rho=\delta_0\}$.  Lemma~\ref{lem:comparison} therefore yields
\[
        U(t,x)\le W(x)+\varepsilon_\tau,
        \qquad \tau\le t\le S,
        \quad 0\le\rho(x)\le\delta_0.
\]
Letting $\tau\downarrow0$ and using continuity gives
\[
        0\le U(t,x)\le W(x),
        \qquad 0\le t\le S,
        \quad 0\le\rho(x)\le\delta_0.
\]
On $\partial\Omega$ the tangential derivatives of $U$ vanish because $u(t,x)=0$ for all $x\in\partial\Omega$ and positive times.  The preceding inequality and $U\ge0$ imply
\[
        0\le \partial_n U(t,x)\le \partial_n W(x)=M\kappa,
        \qquad x\in\partial\Omega,
        \quad 0<t<S.
\]
Therefore
\begin{equation}\label{eq:boundary-gradient-bound}
        |\nabla U(t,x)|\le M\kappa,
        \qquad x\in\partial\Omega,
        \quad 0<t<S.
\end{equation}

It remains to exclude a larger interior maximum.  Put
\[
        w=\frac12|\nabla U|^2.
\]

Fix $\tau>0$. By Lemma~\ref{lem:interior-bootstrap}, on every compact
positive-time interior cylinder the solution has the spatial derivatives
needed to differentiate the equation once and to evaluate the resulting
identity pointwise. This is sufficient because the spatial boundary values
of $|\nabla U|$ have already been controlled by
\eqref{eq:boundary-gradient-bound}. Hence, any larger maximum on
\[
[\tau,S']\times\overline{\Omega},
\qquad S'<S,
\]
is attained at a spatially interior point, where the interior regularity
result applies. If such a maximum lies on the terminal slice $t=S'$, the
time derivative is interpreted from the left; at a terminal-time maximum
one still has $(Y_\delta)_t\geq0$. The calculation is therefore classical
on each such finite time cylinder, and the final inequality depends only
on the structural constants displayed below, rather than on any higher
norm of $U$. At any interior point where $\nabla w=0$, differentiating \eqref{eq:tailored-gradient-pde}, multiplying by $U_k$, and summing over $k$ gives
\[
\begin{aligned}
        w_t-a^{ij}(x,\nabla U)w_{ij}
        &=U_k\,\partial_{x_k}a^{ij}(x,\nabla U)U_{ij}
        +U_k\,\partial_{p_\ell}a^{ij}(x,\nabla U)U_{\ell k}U_{ij} \\
        &\quad+U_k\,\partial_{x_k}b^i(x)U_i
        +U_k b^i(x)U_{ik}
        +U_kF_{x_k}(x,U)+F_r(x,U)|\nabla U|^2 \\
        &\quad-a^{ij}(x,\nabla U)U_{ki}U_{kj}.
\end{aligned}
\]
The identities $w_\ell=U_kU_{k\ell}=0$ imply, for every $\ell$ and every $i$, that
\[
        U_kU_{k\ell}=0,
        \qquad
        U_kU_{ki}=0.
\]
Consequently the term containing $\partial_{p_\ell}a^{ij}$ and the transport term $U_kb^iU_{ik}$ vanish at such a point.  At such a point, ellipticity and Young's inequality give
\[
 -a^{ij}U_{ki}U_{kj}\le-\nupar|D^2U|^2,
\]
while
\[
 |U_k\partial_{x_k}a^{ij}U_{ij}|
 \le C|\nabla U||D^2U|
 \le \frac{\nupar}{4}|D^2U|^2+C_\nu|\nabla U|^2.
\]
Furthermore,
\[
 |U_k(\partial_{x_k}b^i)U_i|\le C|\nabla U|^2,
 \qquad
 |U_kF_{x_k}(x,U)|\le C|\nabla U|\le C(1+w),
\]
and $F_r|\nabla U|^2\le0$.  Therefore
\[
        w_t-a^{ij}(x,\nabla U)w_{ij}
        \le
        -\frac{3\nupar}{4}|D^2U|^2+C_0(1+w)
        \le C_0(1+w)
\]
at every interior point where $\nabla w=0$.  Here $C_0$ depends only on the data listed in the statement.

Choose $C>C_0$ and define
\[
        Y(t,x)=e^{-Ct}(w(t,x)+1).
\]
For $\delta>0$ set $Y_\delta=Y-\delta(t-\tau)$.  If $Y_\delta$ attained a positive maximum at an interior point with time larger than $\tau$, then $\nabla w=0$ there and
\[
\begin{aligned}
        0&\le (Y_\delta)_t-a^{ij}(x,\nabla U)(Y_\delta)_{ij}\\
        &=e^{-Ct}\bigl(w_t-a^{ij}(x,\nabla U)w_{ij}-C(w+1)\bigr)-\delta<0,
\end{aligned}
\]
a contradiction.  Thus, on $[\tau,S']\times\overline\Omega$, the maximum of $Y_\delta$
lies on the initial slice $t=\tau$ or on the spatial boundary.  Therefore, after letting $\delta\downarrow0$, the maximum of $Y$ on
$[\tau,S']\times\overline{\Omega}$ is attained either on the initial
slice $t=\tau$ or on the spatial boundary.
Since the resulting bound is independent of $S'<S$, we may let $S'\uparrow S$.  Since $U(t)\to U^0$ in $C^1(\overline\Omega)$ as $t\downarrow0$, the contribution from the time slice $t=\tau$ converges to $\frac12\norm{\nabla U^0}_{L^\infty}^2$ as $\tau\downarrow0$.  Combining this with \eqref{eq:boundary-gradient-bound} and then letting $\tau\downarrow0$ gives
\[
        \sup_{(0,S)\times\Omega}w
        \le
        e^{CS}\left(1+\frac12\max\left\{
        \norm{\nabla U^0}_{L^\infty(\Omega)}^2,
        M^2\kappa^2
        \right\}\right)-1.
\]
Since \(w=\frac12|\nabla U|^2\), the desired gradient bound follows.
\end{proof}

\end{document}
